\chaptertitle{第七章\quad 推广与展望}

前六章我们以二维平面为主要舞台，建立了一套"几何变换→矩阵→行列式→特征值→齐次坐标→二次型"的知识链条。这条路走到这里，你已经掌握了线性代数最核心的五六个概念——而且每一个都是从具体的几何图像中生长出来的，而不是从抽象定义开始的。

这一章做两件事。第一，把前六章中限于二维的结论系统地推广到任意维度——你会发现大部分内容"只是多加了几行几列"，逻辑一脉相承。第二，介绍几个前六章没能覆盖、但在后续学习和应用中非常重要的工具——克拉默法则、拉普拉斯展开、奇异值分解——作为你继续探索线性代数的路标。

\sectiontitle{7.1\quad 矩阵的推广——从方阵到任意形状}

\textbf{从$2\times2$到$m\times n$。}

前六章的矩阵大多是$2\times2$的方阵（齐次坐标用了$3\times3$）。但矩阵可以是任意行数$m$和列数$n$的组合。一个$m\times n$矩阵有$m$行$n$列：
\[
A = \begin{pmatrix}
a_{11} & a_{12} & \cdots & a_{1n}\\
a_{21} & a_{22} & \cdots & a_{2n}\\
\vdots & \vdots & \ddots & \vdots\\
a_{m1} & a_{m2} & \cdots & a_{mn}
\end{pmatrix}
\]
矩阵乘法的规则不变——左行乘右列，对应相乘再相加——只是维度需要匹配：$m\times n$矩阵乘$n\times p$矩阵得到$m\times p$矩阵。$m\times n$矩阵乘$n\times1$列向量得到$m\times1$列向量。

\textbf{方阵与非方阵的区别。}

方阵（$m=n$）有行列式、迹、逆矩阵（可逆时）和特征值。非方阵（$m\neq n$）没有行列式、没有逆矩阵、也没有特征值——但这些概念都有推广的版本，我们会在7.4节看到。

\textbf{矩阵的秩——一个核心的数字。}

对于一个$m\times n$矩阵，一个非常重要的量是它的\textbf{秩}（rank）。秩可以直观地理解为：矩阵的列向量中线性无关的向量个数。或者等效地说：矩阵代表的变换把空间映射到多少维的子空间中去。

前六章中我们已经反复见到秩的极端情况：
\begin{itemize}
\item 可逆$2\times2$矩阵的秩为$2$（列向量不共线，映射到整个二维平面）。
\item 投影矩阵的秩为$1$（两个列向量共线，映射退化为一条直线）。
\item 零矩阵的秩为$0$（一切映射到原点）。
\end{itemize}

对于$m\times n$矩阵，秩始终满足：秩 $\leq \min(m,n)$。秩的概念把前面学到的"线性相关/无关"和"降维"统一成了一个数字。在今后的数据分析、机器学习等应用中，秩是判断数据中"有多少独立信息"的核心工具。

\begin{exercise}
\item 判断矩阵$\begin{pmatrix}1&2&3\\2&4&6\end{pmatrix}$的秩。提示：观察三列之间的关系。
\item 一个$3\times5$矩阵的秩最大是多少？最小是多少？
\end{exercise}

有了矩阵的秩这个工具，我们就可以讨论任意形状的矩阵了。但前六章中还有一个只对"方阵"定义的工具——行列式。当矩阵从$2\times2$扩大到$n\times n$时，行列式如何计算？第三章的展开公式只到三阶，更一般的方法是什么？

\sectiontitle{7.2\quad $n$阶行列式与拉普拉斯展开}

\textbf{从二阶、三阶到$n$阶。}

三阶行列式的计算（第三章3.4节）按照"主对角线方向取正、副对角线方向取负"的规则展开。但从四阶开始，这种口诀就不再适用了——需要更系统的方法。

对于一般的$n\times n$矩阵$A=(a_{ij})$，行列式定义为：
\[
\det A = \sum_{\sigma}(-1)^{\operatorname{sgn}(\sigma)}\,a_{1\sigma(1)}a_{2\sigma(2)}\cdots a_{n\sigma(n)}
\]
其中求和遍及$\{1,2,\ldots,n\}$的所有$n!$种排列$\sigma$，$\operatorname{sgn}(\sigma)$为排列的符号。这个定义虽然严格，但直接计算需要$n!$项——对于$n=10$，这已经是$3\,628\,800$项——在计算上完全不实用。

\textbf{拉普拉斯展开——逐行降阶。}

实际计算$n$阶行列式的最常用方法之一是\textbf{拉普拉斯展开}（也称按行/按列展开）。对于任意$n\times n$矩阵，可以按第$i$行展开：
\[
\det A = \sum_{j=1}^{n} (-1)^{i+j}\,a_{ij}\,M_{ij}
\]
其中$M_{ij}$是划去第$i$行和第$j$列后剩下的$(n-1)\times(n-1)$子矩阵的行列式，称为$a_{ij}$的\textbf{余子式}。$(-1)^{i+j}M_{ij}$称为$a_{ij}$的\textbf{代数余子式}。

这个公式将$n$阶行列式转化为$n$个$(n-1)$阶行列式的带符号和。递归应用，最终回到我们熟悉的二阶行列式$ad-bc$。

\begin{example}
用拉普拉斯展开计算四阶行列式：$\begin{vmatrix}2&0&1&0\\0&3&0&2\\1&0&2&0\\0&4&0&3\end{vmatrix}$。
\end{example}
\begin{solution}
按第一行展开（第一行有两个非零元，计算量较小）：
\begin{align*}
\det &= 2\cdot(-1)^{1+1}\begin{vmatrix}3&0&2\\0&2&0\\4&0&3\end{vmatrix}
+ 0 + 1\cdot(-1)^{1+3}\begin{vmatrix}0&3&2\\1&0&0\\0&4&3\end{vmatrix}
+ 0\\[2pt]
&= 2\cdot\begin{vmatrix}3&0&2\\0&2&0\\4&0&3\end{vmatrix}
+ 1\cdot\begin{vmatrix}0&3&2\\1&0&0\\0&4&3\end{vmatrix}
\end{align*}
第一个三阶行列式（按第二行展开，因为该行只有一个非零元）：
\[
\begin{vmatrix}3&0&2\\0&2&0\\4&0&3\end{vmatrix}
= 2\cdot(-1)^{2+2}\begin{vmatrix}3&2\\4&3\end{vmatrix}
= 2\cdot(9-8)=2
\]
第二个三阶行列式（按第二行展开）：
\[
\begin{vmatrix}0&3&2\\1&0&0\\0&4&3\end{vmatrix}
= 1\cdot(-1)^{2+1}\begin{vmatrix}3&2\\4&3\end{vmatrix}
= -1\cdot(9-8)=-1
\]
因此$\det = 2\times2 + 1\times(-1) = 3$。
\end{solution}

计算技巧：选零元最多的行（或列）展开，以最小化计算量。在实际应用中，对较大的$n$，更高效的方法是利用第三章3.5节的性质，将矩阵通过行变换化为上三角形式，行列式即对角元乘积——这正是计算机上大多数行列式算法的核心思想。

\begin{exercise}
\item 用拉普拉斯展开（按第一列）计算$\begin{vmatrix}1&2&0\\3&0&1\\0&2&4\end{vmatrix}$，验证结果是否与按第三章的展开公式一致。
\item 为什么选零元最多的行展开可以减小计算量？
\end{exercise}

拉普拉斯展开让我们能计算任意阶的行列式。而行列式的一个重要应用，就是直接写出线性方程组的解——无需消元，一公式到位。这就是古典的克拉默法则。

\sectiontitle{7.3\quad 克拉默法则}

\textbf{优雅但昂贵的解法。}

对于$n$个未知数$n$个方程的线性方程组$A\vec x = \vec b$，如果$\det A\neq0$（矩阵可逆），克拉默法则给出了每个未知数的一个显式公式：
\[
x_i = \frac{\det A_i}{\det A}
\]
其中$A_i$是把$A$的第$i$列替换为$\vec b$后得到的矩阵。

\begin{example}
用克拉默法则解$\begin{cases}2x+3y=7\\x-2y=-4\end{cases}$。
\end{example}
\begin{solution}
$A=\begin{pmatrix}2&3\\1&-2\end{pmatrix}$，$\det A = 2\times(-2)-3\times1 = -7$。

求$x$：将第一列替换为$\vec b=\begin{pmatrix}7\\-4\end{pmatrix}$，得$A_1=\begin{pmatrix}7&3\\-4&-2\end{pmatrix}$，
$\det A_1 = 7\times(-2)-3\times(-4) = -14+12 = -2$。
$x = \frac{-2}{-7} = \frac{2}{7}$。

求$y$：将第二列替换为$\vec b$，得$A_2=\begin{pmatrix}2&7\\1&-4\end{pmatrix}$，
$\det A_2 = 2\times(-4)-7\times1 = -8-7 = -15$。
$y = \frac{-15}{-7} = \frac{15}{7}$。

这恰好与第二章2.4节中用高斯消元法得到的结果一致。
\end{solution}

\textbf{克拉默法则的适用场合。}

克拉默法则在理论上非常优美——它把解线性方程组归结为行列式的计算。但在实际计算中，直接算$n+1$个$n$阶行列式（每个需要$n!$量级的计算）对于$n>3$就不可接受了。高斯消元法（大约$n^3$量级）要高效得多。

克拉默法则的价值主要在理论推导中——比如在证明某些关于解的连续依赖性或者分析含参数的方程组时，行列式的表达式比消元法的中间步骤更易于处理。

\begin{exercise}
\item 用克拉默法则解$\begin{cases}3x-2y=8\\x+4y=2\end{cases}$，验证与高斯消元法的结果一致。
\item 对于$n=100$的方程组，克拉默法则和高斯消元各需要大约多少次运算？据此说明为什么实际计算中几乎不使用克拉默法则。
\end{exercise}

克拉默法则优雅地连接了"行列式"和"解方程"两个主题。但它和特征值分解有一个共同局限：都只适用于方阵。现实世界的数据矩阵几乎都不是方阵——那么，我们如何对一个任意形状的矩阵提取它的"核心信息"？答案是奇异值分解——特征值的"全能"推广。

\sectiontitle{7.4\quad 奇异值分解简介}

\textbf{方阵特征值的局限性。}

第四章学习的特征值分解$A = PDP^{-1}$非常强大，但它有两个限制：第一，只适用于方阵；第二，不是所有方阵都能对角化。现实世界的数据矩阵几乎从不是方阵——例如$100$个学生每人考$5$门课的成绩表就是一个$100\times5$矩阵。

\textbf{奇异值：特征值的"全能"推广。}

对于\textbf{任意}$m\times n$矩阵$A$（无论是否方阵，无论秩为多少），总存在一个分解：
\[
A = U\Sigma V^{\text T}
\]
其中：
\begin{itemize}
\item $U$是$m\times m$正交矩阵（列向量为$AA^{\text T}$的特征向量）。
\item $V$是$n\times n$正交矩阵（列向量为$A^{\text T}A$的特征向量）。
\item $\Sigma$是$m\times n$的"对角"矩阵——仅在$\Sigma_{11},\Sigma_{22},\ldots,\Sigma_{rr}$（$r$为$A$的秩）的位置上有正数，其余全为$0$。这些正数$\sigma_1\geq\sigma_2\geq\cdots\geq\sigma_r>0$就是$A$的\textbf{奇异值}。
\end{itemize}

\textbf{几何意义。}

奇异值分解的几何解释非常优美——而且是第二章"变换的复合"思想的直接延续。任何$m\times n$矩阵$A$代表的线性变换，总可以分解为三步：
\[
V^{\text T}\text{（旋转）}\quad\longrightarrow\quad
\Sigma\text{（沿坐标轴伸缩，伸缩比就是奇异值）}\quad\longrightarrow\quad
U\text{（旋转）}
\]
这恰好是第六章中处理二次型的"旋转→标准化→再旋转"思路在任意维度、任意形状矩阵上的推广。对比一下：
\begin{itemize}
\item 二次型：$M = P D P^{\text T}$——找到主轴方向，沿主轴伸缩。
\item 奇异值分解：$A = U\Sigma V^{\text T}$——先旋转到正确的方向，沿坐标轴伸缩，再旋转回原位。
\end{itemize}

最大的奇异值$\sigma_1$就是矩阵$A$在"最佳方向"上的放大倍率。所有奇异值的平方$\sigma_i^2$恰好是$A^{\text T}A$的特征值。

\textbf{用途。}

奇异值分解在数据压缩（只保留最大的几个奇异值及对应的向量，丢掉小的）、图像处理、推荐系统（如"你可能会喜欢"的算法）、主成分分析（PCA）等众多领域都是核心工具。不过这些内容已经超出了本书的范围——这里只是给你一个路标，告诉你线性代数走到深处是什么图景。

\begin{exercise}
\item 验证：对$A=\begin{pmatrix}1&0\\0&2\\0&0\end{pmatrix}$（一个$3\times2$矩阵），计算$A^{\text T}A$的特征值。这些特征值与$A$的"放大能力"有什么关系？
\item 思考：如果一个$m\times n$矩阵的秩为$r$，它有多少个非零奇异值？为什么？
\end{exercise}

\sectiontitle{7.5\quad $n$维仿射变换——齐次坐标的推广}

第五章我们学会了用$3\times3$齐次坐标矩阵统一处理二维的线性变换和平移。这套方法在任意维度都适用。

\textbf{$n$维齐次坐标。} 在$n$维空间中，一个点用$n$个坐标$(x_1,\ldots,x_n)$表示。给它多加一维$(x_1,\ldots,x_n,1)$，就可以用$(n+1)\times(n+1)$矩阵同时表达线性变换和平移：
\[
\begin{pmatrix}
a_{11}&\cdots&a_{1n}&t_1\\
\vdots&\ddots&\vdots&\vdots\\
a_{n1}&\cdots&a_{nn}&t_n\\
0&\cdots&0&1
\end{pmatrix}
\]
左上角的$n\times n$子矩阵是线性部分，最后一列的前$n$个分量是平移部分。复合、求逆——一切操作和二维时完全一致，只是数字变多了。

\textbf{应用。} $3\times3$齐次矩阵是全部计算机二维图形的数学基础——你看到的每一个图形界面（按钮、窗口、动画）的平移、旋转和缩放，底层都是$3\times3$矩阵乘法。三维游戏和动画则使用$4\times4$齐次矩阵。这套语言在机器人学（描述机械臂的运动）、计算机视觉（描述摄像机的位姿）中无处不在——而它们的核心思想，你已经全部掌握了。

\begin{exercise}
\item 写出三维空间中先绕$z$轴旋转$90^\circ$再沿$x$轴平移$5$个单位的$4\times4$齐次矩阵。绕$z$轴旋转$90^\circ$的$3\times3$部分为$\begin{pmatrix}0&-1&0\\1&0&0\\0&0&1\end{pmatrix}$。
\item 为什么计算机图形学广泛使用齐次坐标而不是分别处理线性变换和平移？
\end{exercise}

\sectiontitle{7.6\quad $n$维二次型与正定性}

第六章用特征值分类二维圆锥曲线——两个特征值同号是椭圆，一正一负是双曲线，有零是抛物线。这个思路在$n$维中有一个极为重要的推广：\textbf{正定性}。

对于一个$n\times n$对称矩阵$M$，它对应的$n$元二次型为：
\[
Q(\vec x) = \vec x^{\text T} M \vec x
\]
根据$M$的特征值，$Q$被分为三类：
\begin{itemize}
\item 所有特征值$>0$：$Q$称为\textbf{正定}二次型。对任何非零$\vec x$，$Q(\vec x)>0$。几何上，$Q(\vec x)=1$的图像是一个$n$维椭球（二维椭圆的推广）。
\item 所有特征值$<0$：$Q$称为\textbf{负定}二次型。$Q(\vec x)<0$始终成立。
\item 特征值有正有负：$Q$称为\textbf{不定}二次型。$Q(\vec x)$可正可负，图像是$n$维双曲面（二维双曲线的推广）。
\item 有特征值$=0$：$Q$称为\textbf{半定}，图像退化为柱面或更低维的曲面。
\end{itemize}

正定性在应用数学中极其重要。在优化问题中，一个函数的"二阶导数矩阵"（Hessian矩阵）的正定性决定了极值点是极小值、极大值还是鞍点。在数据科学中，协方差矩阵必然是半正定的——这保证了你计算的"方差"永远不会是负数。所有这些应用，本质上都是在问同一个问题：这个对称矩阵的特征值是正还是负？——和你判断一个二维曲线是椭圆还是双曲线完全是一回事。

\begin{exercise}
\item 对称矩阵$M=\begin{pmatrix}2&1&0\\1&2&1\\0&1&2\end{pmatrix}$对应一个三元二次型。凭直觉（不用计算特征值），你认为它的三个特征值的符号可能是什么？$Q(\vec x)=\vec x^{\text T}M\vec x$的符号能确定吗？
\item 对比第六章的圆锥曲线分类表，说明$n$维正定性的概念是二维分类的何种推广。
\end{exercise}

\sectiontitle{7.7\quad 全书回顾与综合练习}

从第一章翻开这本书到现在，你走过了这样一条路：

\begin{center}
\begin{tabular}{c|l|l}
\textbf{章节} & \textbf{核心问题} & \textbf{核心工具}\\\hline
第一章 & 点被变换后去了哪里？ & 旋转、反射、伸缩、投影、切变的坐标公式\\
第二章 & 怎么紧凑地表示和复合变换？ & 矩阵、矩阵乘法、逆矩阵、高斯消元\\
第三章 & 变换把面积/体积放大了几倍？ & 行列式、线性相关与无关、降维\\
第四章 & 哪些方向在变换中保持不变？ & 特征值、特征向量、施密特正交化、对角化\\
第五章 & 平移怎么融入矩阵框架？ & 齐次坐标、仿射变换\\
第六章 & 怎么用矩阵优雅地分类圆锥曲线？ & 二次型、特征值分类法、主轴方向\\
\end{tabular}
\end{center}

这条路的核心思想只有一个：\textbf{用代数的结构去理解几何的动作，用几何的直觉去赋予代数以意义。}

下面是全书范围的综合练习。这些题目跨越多个章节，需要你调动前面学过的各种工具来综合解决。

\subsection*{综合练习}
\begin{exercise}
\item 某变换的矩阵为$A=\begin{pmatrix}2&1\\1&2\end{pmatrix}$。\\
(a) 这个变换把面积放大了几倍？\\
(b) 求出它的特征值和特征向量，并解释特征向量的几何含义。\\
(c) 写出二次型$(x,y)A\begin{pmatrix}x\\y\end{pmatrix}$的表达式，判断其对应的曲线类型。\\
(d) 对这个矩阵做对角化：求$P$和$D$使得$P^{-1}AP=D$。
\item 已知平面上一个变换把$(1,0)$变成$(2,3)$，把$(0,1)$变成$(1,4)$。\\
(a) 写出这个变换的$2\times2$矩阵。\\
(b) 计算它的行列式，说明面积如何变化。\\
(c) 这个变换是否可逆？如果可逆，求逆矩阵。\\
(d) 求这个变换的特征值，并判断是否有特征向量方向不变。
\item 考虑二次曲线$2x^2+2xy+2y^2=3$。\\
(a) 写出对应的对称矩阵$M$。\\
(b) 求$M$的特征值，判断曲线类型。\\
(c) 求主轴方向。\\
(d) 写出旋转后的标准方程。
\item 矩阵$A=\begin{pmatrix}1&1&0\\0&1&1\\0&0&1\end{pmatrix}$表示三维空间中的一个变换。\\
(a) 判断三个列向量是否线性无关。\\
(b) 计算行列式，解释几何意义（提示：这是切变，体积不变）。\\
(c) 求$A$的特征值。$A$可以对角化吗？
\item 试用自己的语言解释：为什么行列式为$0$的矩阵没有逆矩阵？你的解释需要同时包含几何角度（变换对空间做了什么）和代数角度（方程$A\vec x=\vec b$的解的情况）。
\item 回顾全书的六个核心概念，选择其中两个，用一个具体的$2\times2$矩阵为例，展示这两个概念之间的关系（例如：行列式与特征值的关系，矩阵乘法与变换复合的关系，对角化与二次型的关系等）。
\end{exercise}

\sectiontitle{课后练习}

\subsection*{A组\quad 基础巩固}
\begin{exercise}
\item 求矩阵$\begin{pmatrix}1&2&3\\0&1&2\\0&0&1\end{pmatrix}$的逆矩阵。提示：用高斯消元法或观察其作为切变矩阵的逆。
\item 用拉普拉斯展开计算$\begin{vmatrix}2&1&0&0\\1&2&1&0\\0&1&2&1\\0&0&1&2\end{vmatrix}$。
\item 用克拉默法则解$\begin{cases}2x+y=5\\x-3y=-1\end{cases}$。
\item 判断：一个$3\times4$矩阵是否有行列式？是否有特征值？其秩最大为多少？
\item 若$A$是一个$2\times2$矩阵，其特征值为$3$和$-2$。求$A^{\text T}A$的特征值。
\item 简述克拉默法则和高斯消元法各自的优缺点。
\end{exercise}

\subsection*{B组\quad 挑战提升}
\begin{exercise}
\item 证明：对于任意$m\times n$矩阵$A$，$A^{\text T}A$总是对称矩阵，且其特征值均非负。
\item 设$A$是$n$阶上三角矩阵（主对角线以下全为零）。证明$\det A$等于主对角线元素的乘积。提示：反复按最后一行做拉普拉斯展开。
\item 一个$2\times2$矩阵的奇异值分解中，$\Sigma=\begin{pmatrix}3&0\\0&1\end{pmatrix}$，$V=I$（单位矩阵），$U=\begin{pmatrix}0&1\\1&0\end{pmatrix}$。求原矩阵$A$。这个矩阵在几何上做了什么？
\item 用拉普拉斯展开证明范德蒙德行式：
$\begin{vmatrix}1&x_1&x_1^2\\1&x_2&x_2^2\\1&x_3&x_3^2\end{vmatrix}=(x_2-x_1)(x_3-x_1)(x_3-x_2)$。
\item 用克拉默法则解$\begin{cases}2x-3y+z=1\\x+y-2z=4\\3x-y+z=2\end{cases}$，写出$A_1,A_2,A_3$并分别计算行列式。
\item 一个$3\times3$对称矩阵的特征值为$4,1,-2$。写出它对应的二次型经过正交变换后的标准形，并判断该二次型的正定性。
\item 利用拉普拉斯展开和行列式性质计算$\begin{vmatrix}1&2&0&1\\2&0&1&0\\0&1&3&2\\1&0&2&1\end{vmatrix}$。
\end{exercise}

\newpage
